\section*{Hoofdstuk \ref{ch:b3:x-ray-diffraction} --- Kristallografie en röntgendiffractie}

\begin{solution}{exo:b3:x-ray-diffraction:1}
$d = a/\sqrt N$: $d_{111} = \qty{325.66}{pm}$, $d_{200} = \qty{282.03}{pm}$, $d_{220} =
\qty{199.42}{pm}$.
\end{solution}

\begin{solution}{exo:b3:x-ray-diffraction:2}
Kubisch vlakgecentreerd: de eerste reflecties zijn (111) en (200). $\sin\theta = \lambda\sqrt N/2a$: $2\theta =
\ang{43.32}$ en \ang{50.45}.
\end{solution}

\begin{solution}{exo:b3:x-ray-diffraction:3}
P: alle zeven. I ($h + k + l$ even): 110, 200, 211, 220. F (ongemengd): 111, 200, 220.
\end{solution}

\begin{solution}{exo:b3:x-ray-diffraction:4}
$\mathbf a^* = \mathbf e_x/a$, $\mathbf b^* = \mathbf e_y/b$: een rechthoekig
rooster met zijden $1/a$
en $1/b$, waarbij de lange zijde van het directe rooster de korte wordt.
\end{solution}

\begin{solution}{exo:b3:x-ray-diffraction:5}
Verhoudingen van $\sin^2\theta$: $1 : 2 : 3 : 4 : 5$, dus $N = 2, 4, 6, 8, 10$: ruimtegecentreerd (110,
200, 211, 220, 310). Uit de eerste lijn is $d_{110} = \lambda/2\sin(\ang{20.135}) =
\qty{223.77}{pm}$ en $a = d\sqrt2 = \qty{316.5}{pm}$: wolfraam.
\end{solution}

\begin{solution}{exo:b3:x-ray-diffraction:6}
$Z = \rho N_Aa^3/M = 1.99 \times 6.022\times10^{23} \times (6.2929\times10^{-8})^3/74.6 = 4.0$.
\end{solution}

\begin{solution}{exo:b3:x-ray-diffraction:7}
$\beta = \qty{0.40}{\degree} = \qty{6.98e-3}{rad}$, $\cos\theta = \cos\ang{15.85} = 0.962$:
$L = 0.9 \times 0.15406/(6.98\times10^{-3} \times 0.962) = \qty{21}{nm}$.
\end{solution}

\begin{solution}{exo:b3:x-ray-diffraction:8}
$F_{hkl} = f_{\ce{Cs+}} + f_{\ce{Cl-}}(-1)^{h+k+l}$: $F_{100} = 54 - 18 = 36$, $F_{110} = 72$:
intensiteiten in de verhouding 1 : 4. Het centrum wordt bezet door een
ander ion dan de hoekpunten: het rooster is primitief (een
ruimtegecentreerd rooster vereist een identieke inhoud op beide punten), en
100 is aanwezig.
\end{solution}

\begin{solution}{exo:b3:x-ray-diffraction:9}
Elk atoom heeft een kopie verschoven over $(\frac12,\frac12,0)$, wat $F$ vermenigvuldigt met $1 +
\eu^{\iu\pi(h+k)} = 1 + (-1)^{h+k}$, nul voor oneven $h + k$.
\end{solution}

\begin{solution}{exo:b3:x-ray-diffraction:10}
Quasikristallen hebben orde op lange afstand zonder periodiciteit: ze
hebben geen rooster, zodat de beperking, die over roosters gaat, niet op
hen van toepassing is. Hun scherpe diffractievlekken leidden ertoe een
kristal te definiëren als elke vaste stof met een in wezen discreet
diffractiepatroon, periodiek of niet.
\end{solution}

\begin{solution}{exo:b3:x-ray-diffraction:11}
De $c$-glijspiegeling beeldt $(x, y, z)$ af op $(x, \frac12 - y, \frac12 + z)$. Voor $k = 0$ zijn de twee
fasefactoren $\eu^{2\pi\iu(hx + lz)}$ en $\eu^{2\pi\iu(hx + lz + l/2)} = (-1)^l
\eu^{2\pi\iu(hx+lz)}$: hun som is nul voor oneven $l$.
\end{solution}

\begin{solution}{exo:b3:x-ray-diffraction:12}
Een inversie, een spiegeling of een glijspiegeling zet een molecuul om in
zijn spiegelbeeld, zodat een kristal dat er een bevat, beide enantiomeren
zou bevatten. Een zuiver enantiomeer kristalliseert alleen in de 65
\omterm{def:b3:x-ray-diffraction:space-group}{ruimtegroepen} die opgebouwd zijn uit rotaties, schroefassen en translaties
(sohnckegroepen), zoals \textit{P}$2_12_12_1$. Volgens de wet van Friedel is
het patroon van het ene enantiomeer gelijk aan dat van het andere; de
absolute configuratie wordt verkregen uit de kleine afwijkingen die
veroorzaakt worden door anomale verstrooiing aan zwaardere atomen.
\end{solution}

\begin{solution}{pb:b3:x-ray-diffraction:1}
\textbf{1.} $\theta = \ang{14.171}$, $d = \lambda/2\sin\theta = \qty{314.64}{pm}$.
\textbf{2.} 314.64, 222.49, 181.66, 157.32, 140.71, \qty{128.45}{pm}.
\textbf{3.} 1.000, 2.000, 3.000, 4.000, 5.000, 6.000.
\textbf{4.} 1, 2, 3, 4, 5, 6.
\textbf{5.} Ja, als (100), (110), (111), (200), (210), (211) van een
primitieve kubische cel met $a' = \qty{314.6}{pm}$.
\textbf{6.} $N = 7$ (en 15, 23\dots), dat geen som van drie kwadraten is;
het eerste verschil zou bij de achtste lijn optreden. Zes lijnen kunnen een
P-cel met ribbe $a'$ niet onderscheiden van een F-cel met ribbe $2a'$.
\textbf{7.} $N = 4, 8, 12, 16, 20, 24$: (200), (220), (222), (400), (420), (422).
\textbf{8.} $a = 2d_{200} = \qty{629.3}{pm}$.
\textbf{9.} \ce{KCl} (\qty{629.29}{pm}).
\textbf{10.} $F = 4[f(\ce{K+}) - f(\ce{Cl-})]$ voor indices die alle oneven
zijn, en de twee ionen hebben elk 18
elektronen: de amplitudes heffen elkaar op.
\textbf{11.} Ja voor beide: \ce{Na+} (10) en \ce{Cl-} (18), \ce{K+} (18) en
\ce{Br-} (36) verstrooien verschillend; (111) van \ce{NaCl} zou bij
\ang{27.36} liggen, die van \ce{KBr} bij \ang{23.38}, beide binnen het
opgenomen bereik.
\textbf{12.} Röntgenstralen zien \ce{K+} en \ce{Cl-} als vrijwel identiek;
identieke ionen in een steenzoutschikking vormen een eenvoudig kubisch
rooster met ribbe $a/2$, een schijnbare cel.
\textbf{13.} Vier \ce{KCl}.
\textbf{14.} $\rho = 4 \times 74.6/(6.022\times10^{23} \times (6.293\times10^{-8})^3) =
\qty{1.99}{g/cm^3}$.
\textbf{15.} Zes en zes (octaëders van het andere ion).
\textbf{16.} $a/2 = \qty{314.6}{pm}$.
\textbf{17.} (440), $N = 32$, bij \ang{87.65} (omdat (333)/(511), met alle
indices oneven, ontbreken).
\textbf{18.} Intensiteiten veranderen met de voorkeursoriëntatie van de
kristallieten, de absorptie, de afname van de verstrooiingsfactoren en de
thermische beweging; posities hangen alleen af van het rooster.
\textbf{19.} Uit $d = \lambda/2\sin\theta$ volgt $\Delta d/d = -\cot\theta\,\Delta\theta$: dezelfde
hoekfout geeft bij grote $\theta$ een kleinere relatieve fout.
\textbf{20.} $d_{420} = \lambda/2\sin\ang{33.190} = \qty{140.71}{pm}$, $a = d\sqrt{20} =
\qty{629.3}{pm}$.
\textbf{21.} $L = 0.9 \times 0.15406/(4.36\times10^{-3} \times \cos\ang{33.19}) = \qty{38}{nm}$.
\textbf{22.} Nee: $\beta = \ang{0.0095}$, ver onder een typische
instrumentele breedte.
\textbf{23.} \textbf{$a = \qty{629.3}{pm}$: het poeder is kaliumchloride},
waarvan de oneven reflecties tot zwijgen gebracht worden door twee ionen
met hetzelfde aantal elektronen.
\end{solution}
